% Part: methods
% Chapter: proofs
% Section: introduction

\documentclass[../../../include/open-logic-section]{subfiles}

\begin{document}

\olfileid{mth}{prf}{int}

\olsection{પરિચય}

પ્રારંભિક તર્કશાસ્ત્રના તમારા અનુભવને લીધે તમે કદાચ કોઈ
!!a{derivation} તંત્રથી પરિચિત હશો---સંભવતઃ સ્વાભાવિક નિષ્કર્ષણનું
અથવા ફિચ-શૈલીનું !!{derivation} તંત્ર, કે પછી સાબિતી-વૃક્ષનું તંત્ર.
આવા તંત્રોમાં તમે કદાચ સાબિતીઓ કરી હશે: !!a{formula} સાબિત કર્યું
હશે અથવા આપેલી દલીલ પ્રામાણ્ય ધરાવે છે તે બતાવ્યું હશે. એ માટે તમે
ઇચ્છિત અંતિમ પરિણામ મળે ત્યાં સુધી તંત્રના નિયમો લાગુ પાડ્યા હતા.
પરંતુ તર્કશાસ્ત્ર \emph{વિશે} વિચારતી વખતે પણ આપણે બાબતો સાબિત કરીએ
છીએ અને મોટા ભાગે !!a{derivation} તંત્ર વાપરતા નથી. હકીકતમાં, આપણે
વિચારીએ છીએ એવી મોટા ભાગની સાબિતીઓ સંપૂર્ણપણે પ્રથમ-ક્રમના તર્કની
ભાષામાં નહીં, પરંતુ અંગ્રેજી જેવી સામાન્ય ભાષામાં (અહીં ગુજરાતીમાં;
ક્યારેક સાંકેતિક ભાષા સાથે) લખાય છે. આવી સાબિતી રચતી વખતે શરૂઆતમાં
મૂંઝવણ થાય: !!a{derivation} તંત્ર વિના કંઈક કેવી રીતે સાબિત કરું?
શરૂઆત ક્યાંથી કરું? મારી સાબિતી સાચી છે તે કેવી રીતે જાણું?

સાબિતી કરવાનો પ્રયાસ કરતાં પહેલાં સાબિતી શું છે અને તે કેવી રીતે
રચાય છે તે જાણવું જરૂરી છે. નામ સૂચવે છે તેમ, \emph{સાબિતી}નો
હેતુ કોઈ બાબત સાચી છે તે બતાવવાનો છે. તેને સંવાદ તરીકે વિચારો:
કોઈ પૂછે કે કોઈ બાબત સાચી છે કે નહીં, દાખલા તરીકે, બે સિવાયની
દરેક અવિભાજ્ય સંખ્યા એકી છે કે નહીં. માત્ર ``હા'' કહેવું પૂરતું
નથી; તે \emph{શા માટે} એવું છે તે જાણવા માગી શકે. આ કિસ્સામાં તમે
તેને સાબિતી આપશો.

રોજિંદી ચર્ચામાં ઉત્તરનો માત્ર સંકેત આપવો અથવા અધૂરો ઉત્તર આપવો
પૂરતો હોઈ શકે. પરંતુ તર્કશાસ્ત્ર અને ગણિતમાં આપણે કડક સાબિતી
ઇચ્છીએ છીએ---કોઈ બાબત \emph{જરાય} શંકા રહે નહીં એ રીતે સાચી છે તે
બતાવવા ઇચ્છીએ છીએ. એટલે આપણી સાબિતીનું દરેક પગલું ન્યાયસંગત હોવું
જોઈએ અને તેનો આધાર યોગ્ય હોવો જોઈએ (એટલે, તમે વાપરો છો તે
પૂર્વધારણા ખરેખર તમે સાબિત કરતા પ્રમેયના વિધાનમાં અપાઈ હોવી જોઈએ,
વાપરેલી વ્યાખ્યાઓ યોગ્ય રીતે લાગુ પાડેલી હોવી જોઈએ, અને જે
આધારો આપો છો તે સાચાં અનુમાનો હોવાં જોઈએ, વગેરે).

સામાન્ય રીતે આપણે કોઈ વિધાન સાબિત કરીએ છીએ. સાબિત કરવાના
વિધાનોને જુદાં નામ આપીએ છીએ: પ્રતિજ્ઞાપન, પ્રમેય, ઉપપ્રમેય અથવા
પરિણામ. પ્રતિજ્ઞાપન સાબિત કરવા યોગ્ય મૂળભૂત વિધાન છે: નોંધવા જેટલું
મહત્ત્વનું, પરંતુ કદાચ ખાસ ઊંડું નહીં કે વારંવાર વપરાતું નહીં.
પ્રમેય મહત્ત્વનું પ્રતિજ્ઞાપન છે. તેની સાબિતી ઘણી વાર કેટલાંક
પગલાંમાં વહેંચાય છે અને ક્યારેક તેને સૌપ્રથમ સાબિત કરનાર વ્યક્તિનું
નામ આપવામાં આવે છે (જેમ કે કૅન્ટરનું પ્રમેય, લેવેન્હાઇમ--સ્કોલેમનું
પ્રમેય) અથવા તે જે બાબત અંગે છે તેનું નામ આપવામાં આવે છે (જેમ કે
પૂર્ણતાનું પ્રમેય). ઉપપ્રમેય એ વધુ મહત્ત્વના પરિણામની સાબિતીમાં
વપરાતું પ્રતિજ્ઞાપન અથવા પ્રમેય છે. ગૂંચવણ એ છે કે ક્યારેક ઉપપ્રમેય
પોતે જ મહત્ત્વનું પરિણામ હોય છે અને તેને રજૂ કરનાર વ્યક્તિના નામથી
પણ ઓળખાય છે (જેમ કે ઝોર્નનું ઉપપ્રમેય). પરિણામ એ બીજા કોઈ પરિણામમાંથી
સહેલાઈથી નીકળતું પરિણામ છે.

સાબિત કરવાના વિધાનમાં ઘણી વાર એવી પૂર્વધારણાઓ હોય છે જે સ્પષ્ટ કરે
છે કે આપણે કઈ જાતની વસ્તુઓ વિશે કંઈક સાબિત કરીએ છીએ. તે ``ધારો કે
$!A$ એ $!B \lif !C$ સ્વરૂપનું !!a{formula} છે'' અથવા ``ધારો કે
$\Gamma \Proves !A$'' એવા શબ્દોથી શરૂ થઈ શકે. આ પ્રતિજ્ઞાપન,
પ્રમેય અથવા ઉપપ્રમેયની \emph{પૂર્વધારણાઓ} છે અને તમારી સાબિતીમાં તમે
તેમને સાચી માની શકો. તે આપણે શું સાબિત કરીએ છીએ તે મર્યાદિત કરે
છે અને જે વસ્તુઓ વિશે વાત કરીએ છીએ તેમનાં નામ પણ રજૂ કરે છે.
ઉદાહરણ તરીકે, તમારું પ્રતિજ્ઞાપન ``ધારો કે $!A$ એ $!B \lif !C$
સ્વરૂપનું !!a{formula} છે'' એમ શરૂ થાય, તો તમે ફક્ત ચોક્કસ જાતનાં
સૂત્રો (એટલે શરતી વિધાનો) વિશે કંઈક સાબિત કરો છો; અને તમારી
સાબિતીમાં આવતું $!B \lif !C$ કોઈ પણ મનસ્વી શરતી વિધાન છે એમ
સમજવાનું છે.

\end{document}
